\documentclass[11pt,a4paper]{article}
\usepackage[utf8]{inputenc}
\usepackage[italian]{babel}
\usepackage{geometry}
\geometry{top=2.5cm,bottom=2.5cm,left=2.5cm,right=2.5cm}
\usepackage{amsmath,amssymb}
\usepackage{tcolorbox}
\tcbuselibrary{skins,breakable}
\usepackage{booktabs}
\usepackage{tabularx}
\usepackage{enumitem}
\usepackage{xcolor}
\usepackage{listings}
\usepackage{hyperref}
\hypersetup{
    colorlinks=true,
    linkcolor=blue!70!black,
    citecolor=blue!70!black,
    urlcolor=teal!70!black
}

\lstset{
    basicstyle=\ttfamily\small,
    breaklines=true,
    frame=single,
    backgroundcolor=\color{gray!10},
    keywordstyle=\color{blue}\bfseries,
    commentstyle=\color{green!50!black}\itshape,
    stringstyle=\color{red!70!black}
}

\newtcolorbox{questionbox}[1][]{
    colback=gray!5!white,
    colframe=gray!60!black,
    fonttitle=\bfseries,
    title={#1},
    breakable
}

\newtcolorbox{solutionbox}[1][]{
    colback=blue!5!white,
    colframe=blue!70!black,
    fonttitle=\bfseries,
    title={#1},
    breakable
}

\newtcolorbox{altbox}[1][]{
    colback=orange!5!white,
    colframe=orange!80!black,
    fonttitle=\bfseries,
    title={#1},
    breakable
}

\begin{document}

\title{\textbf{Sicurezza dei Sistemi Informativi} \\ \large Appello d'Esame del 14 Febbraio 2025 \\ \small Anno Accademico 2024-2025}
\author{Politecnico di Torino --- Docente: Prof. Cataldo Basile}
\date{2025-02-14}
\maketitle


\section{Testo Ufficiale dell'Esame}

\begin{questionbox}[Parte 1: Quiz Ufficiale su Piattaforma d'Esame]
Valutare ciascuna delle affermazioni selezionando quelle corrette per ciascun gruppo di quesiti:
\begin{enumerate}[label=\textbf{1.\arabic*}]
    \item \textbf{Sicurezza di Rete e Monitoraggio:}
    \begin{enumerate}[label=\alph*.]
        \item Un WAF, posto in un gateway in modalit\`a stealth, \`e in grado di filtrare il traffico di applicazioni web anche se queste sono cifrate mediante TLS.
        \item In un'architettura di tipo \textit{screened host gateway} si pu\`o mascherare la rete interna solo se il router di bordo implementa anche funzioni di NAT.
        \item Data breach \`e un incidente per cui \`e stato confermato che dei dati sono stati effettivamente rivelati (non solo esposti) a persone non autorizzate.
        \item Definire quale \`e il traffico autorizzato \`e uno dei problemi principali per la corretta configurazione di un firewall.
        \item Gli IDS avanzati basati su AI permettono di identificare e reagire alle minacce in maniera efficace.
        \item Un application-level filter permette di definire regole di autorizzazione pi\`u granulari rispetto ad un packet filter, perch\'e \`e in grado di comprendere la semantica e la sintassi del protocollo che si intende filtrare.
        \item Un network-based IDS monitora il traffico di una sottorete ma non \`e pensato per identificare attacchi provenienti da Internet.
        \item In un SOC non dovrebbero esserci esperti di analisi forense, perch\'e lo scopo principale \`e monitorare, identificare ed evitare gli attacchi.
        \item Un SIEM serve a collezionare e analizzare dati ai fini di monitoring oltre a permettere di generare report utili anche per verifiche di compliance.
    \end{enumerate}

    \item \textbf{Protocolli di Autenticazione e Controllo Accessi di Rete:}
    \begin{enumerate}[label=\alph*.]
        \item PPP serve ad incapsulare pacchetti di livello 3 per trasferimenti punto-punto.
        \item Per come \`e progettato, EAP richiede frequenti cambiamenti allo standard, in pratica ogni volta che \`e necessario aggiungere un nuovo tipo di autenticazione.
        \item Le funzioni di sicurezza intrinseche di RADIUS sono insufficienti.
        \item Il protocollo RADIUS permette di incapsulare messaggi EAP.
        \item I server RADIUS possono fare da proxy, ma solo verso altri server RADIUS.
        \item EAP \`e un middleware che permette di usare diversi tipi di autenticazione con diversi tipi di protocolli di comunicazione di livello 2.
        \item Eduroam richiede una condivisione minima di segreti utenti tra istituzioni federate.
        \item PAP \`e ancora considerato un metodo di autenticazione valido.
        \item Se si usa 802.1X, gli AP non devono necessariamente autenticare gli utenti direttamente.
    \end{enumerate}

    \item \textbf{Identit\`a Digitale e Meccanismi di Autenticazione:}
    \begin{enumerate}[label=\alph*.]
        \item Nel Digital Authentication Model del NIST, il Credential Service Provider ha il compito di emettere o registrare le credenziali degli utenti e gli altri attributi ad essi associati.
        \item Un dictionary attack ha, a parit\`a di tempo a disposizione, una maggiore probabilit\`a di identificare le password degli utenti rispetto ad un approccio a forza bruta, considerando che molti sceglieranno password deboli.
        \item Un password wallet \`e un esempio semplice ed efficace di Single Sign-On che rispetta tutti i requisiti di un SSO multi-dominio e multi-applicazione.
        \item In Kerberos, gli Application Server ricevono dal Ticket Granting Server le chiavi degli utenti.
        \item Il token FIDO2 svolge anche il ruolo di fattore di autenticazione di tipo "possesso".
        \item L'autenticazione mediante dispositivo personale, nella versione che sfrutta l'inerenza, viene ritenuta molto sicura.
        \item Un approccio di tipo Zero-Knowledge Proof include anche il caso in cui un Verifier possa fare centinaia di domande ad un Prover prima di garantire l'autenticazione.
        \item I CRA asimmetrici sono resistenti allo sniffing.
    \end{enumerate}

    \item \textbf{Infrastruttura a Chiave Pubblica (PKI) e Firme Digitali:}
    \begin{enumerate}[label=\alph*.]
        \item L'OCSP mi permette di identificare le firme che non sono valide perch\'e posso capire se il certificato era valido al momento della firma.
        \item La fiducia nei certificati digitali degli utenti, in realt\`a, dipende dalla fiducia nelle autorit\`a di certificazione che li hanno emessi (se la crittografia \`e forte ed usata correttamente).
        \item L'algoritmo di Grover render\`a facilmente attaccabili le firme digitali calcolate con gli algoritmi della crittografia classica (es. RSA, DSA).
        \item Un documento contenente una macro pu\`o essere firmato, ma viene garantita l'integrit\`a e autenticazione del codice della macro stessa, non il risultato di quello che verr\`a calcolato o cambiato sul documento.
        \item Un certificato X.509 di tipo Diffie-Hellman riporter\`a tre parametri: il numero primo $p$, il generatore $g$ e $g^x \bmod p$, dove $x$ \`e il parametro privato del possessore del certificato.
        \item Un documento firmato include un Time-Stamp ottenuto da una TSA fidata e riporta come data e ora il 14/02/2025 alle 14:00. Posso quindi escludere che il documento esistesse gi\`a il 13/02/2025.
        \item Dato un documento che deve essere firmato da $n$ firmatari, se si scegliesse di usare firme sequenziali la verifica richiederebbe pi\`u tempo rispetto al caso in cui si scegliesse l'uso di firme parallele.
    \end{enumerate}

    \item \textbf{Sicurezza Software, Standard e Crittografia Pratica:}
    \begin{enumerate}[label=\alph*.]
        \item Una vulnerabilit\`a del tipo buffer overflow potrebbe permettere agli attaccanti di prendere il controllo di un server remoto.
        \item La certificazione di sicurezza aumenta la nostra fiducia nel prodotto certificato perch\'e sappiamo che qualcuno ha verificato che determinate propriet\`a di sicurezza siano soddisfatte.
        \item Il processo di valutazione definito nello standard ITSEC si concentra su efficacia delle soluzioni e correttezza dell'implementazione.
        \item Bisogna cambiare spesso le chiavi crittografiche simmetriche usate per garantire la confidenzialit\`a del traffico di rete, perch\'e gli attaccanti potrebbero ricavare dal traffico cifrato informazioni utili a decifrarlo.
        \item Senza l'opzione \texttt{-no-pad} (OpenSSL), cifrare 16 byte con AES128-ECB \`e preferibile rispetto ad usare AES128-CBC, considerando sicurezza e spreco di dati.
        \item La crittografia a curve ellittiche costruisce algoritmi di crittografia asimmetrica a partire da problemi difficili su campi matematici ottenuti da equazioni del tipo $y^2 = x^3 + ax + b$.
        \item Gli algoritmi di Hash vengono usati per generare le chiavi che garantiscono confidenzialit\`a e integrit\`a+autenticazione nel TLS.
        \item Se l'ultimo blocco di dati cifrato con DES termina con \texttt{0x1324354603030303} verr\`a accettato correttamente sia con che senza l'opzione \texttt{-no-pad} (OpenSSL).
    \end{enumerate}
\end{enumerate}
\end{questionbox}

\begin{questionbox}[Parte 2: Domande a Risposta Aperta (Compito Scritto)]
\textbf{Domanda 1: NIST RMF e Enterprise Patch Management}
\begin{enumerate}[label=\alph*.]
    \item Spiegare cosa \`e il Risk Management Framework (RMF) del NIST (SP 800-37) e quali sono gli obiettivi di tale framework. Elencare e descrivere le 7 fasi e alcune delle operazioni che potrebbero essere svolte in ognuna di esse.
    \item Spiegare cosa \`e l'Enterprise Patch Management (NIST SP 800-40), quali sono i suoi obiettivi operativi e citare almeno due aspetti e sfide rilevanti (es. prioritizzazione basata su CVSS/EPSS, test preventivi di regressione in ambiente di staging).
\end{enumerate}

\textbf{Domanda 2: Progettazione Architettura di Rete Sicura}
Progettare un'infrastruttura di rete sicura per un'azienda dotata di sede centrale, DMZ con servizi esposti e lavoratori remoti. Discutere le tecnologie impiegate (Firewall stateful, VPN IPsec/TLS, RADIUS/802.1X, MFA) e dettagliare la matrice di autorizzazione.
\end{questionbox}

\newpage
\section{Soluzione Dettagliata e Guida Risolutiva}

\begin{solutionbox}[Soluzione Parte 1: Analisi Dettagliata del Quiz Ufficiale]
\textbf{1.1 Sicurezza di Rete e Monitoraggio (Analisi Puntuale delle Opzioni):}
\begin{itemize}
    \item \textbf{a. FALSA}: Un Web Application Firewall (WAF) collocato in modalit\`a \textit{stealth} trasparente (bridge L2) non pu\`o ispezionare n\'e filtrare il traffico applicativo HTTP cifrato con TLS se non possiede i certificati digitali e la relativa chiave privata per terminare o intercettare la sessione crittografica.
    \item \textbf{b. VERA}: Nell'architettura \textit{screened host gateway}, il packet filtering router convoglia il traffico pubblico sul bastion host esposto; per mascherare e proteggere la sottorete interna privata (RFC 1918), il router di bordo deve obbligatoriamente implementare funzioni di Network Address Translation (NAT/NAPT).
    \item \textbf{c. FALSA}: Ai sensi dell'Art. 4.12 del GDPR, per "violazione dei dati personali" (\textit{data breach}) si intende qualsiasi violazione di sicurezza che comporti accidentalmente o in modo illecito la distruzione, la perdita, la modifica, la divulgazione non autorizzata o l'accesso ai dati; l'esposizione indebita di dati non protetti configura gi\`a un breach conclamato anche in assenza di prove di effettiva esfiltrazione attiva.
    \item \textbf{d. VERA}: Definire con precisione cosa costituisca traffico lecito rispetto alle esigenze aziendali (\textit{Authorized Traffic Baseline}) \`e una delle sfide pi\`u critiche ed esposte a fallacie nella configurazione delle regole di un firewall.
    \item \textbf{e. FALSA}: Gli IDS basati su modelli di machine learning e intelligenza artificiale soffrono storicamente di elevati tassi di falsi positivi dovuti al mutamento dei pattern di traffico (\textit{concept drift}); una risposta automatica non supervisionata rischia di bloccare servizi operativi legittimi generando Denial of Service auto-indotti.
    \item \textbf{f. VERA}: Un filtro applicativo (L7) esegue il riassemblaggio dello stack e l'analisi semantica e sintattica del protocollo (es. comandi HTTP, query SQL, intestazioni MIME), consentendo politiche di autorizzazione contestuali inaccessibili a un packet filter L3/L4.
    \item \textbf{g. FALSA}: I Network-based IDS (NIDS) sono concepiti e posizionati strategicamente proprio ai confini perimetrali (a monte o a valle del firewall) per rilevare intrusioni, scansioni ed exploit provenienti da Internet verso le risorse aziendali.
    \item \textbf{h. FALSA}: Il Security Operations Center (SOC) moderno integra necessariamente competenze di Digital Forensics \& Incident Response (DFIR) per raccogliere prove conformi a valore legale, determinare la causa radice (\textit{root cause}) dell'incidente e valutare la persistenza delle minacce.
    \item \textbf{i. VERA}: La piattaforma SIEM (\textit{Security Information and Event Management}) centralizza, normalizza e correla i log di rete, fornendo dashboard di monitoraggio operativo e audit trail conformi a standard di compliance (GDPR, ISO/IEC 27001, PCI-DSS).
\end{itemize}

\textbf{1.2 Protocolli di Autenticazione e Controllo Accessi di Rete:}
\begin{itemize}
    \item \textbf{a. VERA}: Il Point-to-Point Protocol (PPP, RFC 1661) \`e formalmente progettato per incapsulare datagrammi multi-protocollo di livello 3 su collegamenti punto-a-punto diretti.
    \item \textbf{b. FALSA}: L'Extensible Authentication Protocol (EAP, RFC 3748) \`e un framework estensibile per design; supporta l'integrazione di nuovi algoritmi di autenticazione semplicemente assegnando un nuovo \textit{Type Code} numerico, senza richiedere modifiche allo standard di base.
    \item \textbf{c. VERA}: Le specifiche originarie di RADIUS (RFC 2865) sono crittograficamente carenti: proteggono solo il campo password tramite un meccanismo debole basato su MD5 e segreto condiviso, trasmettono tutti gli altri attributi utente in chiaro e non garantiscono integrit\`a end-to-end su tutti i pacchetti (motivo per cui si adotta RadSec con TLS, RFC 6614).
    \item \textbf{d. VERA}: Lo standard RFC 3579 specifica formalmente come i messaggi EAP siano incapsulati all'interno degli attributi \texttt{EAP-Message} e protetti dall'attributo \texttt{Message-Authenticator} di RADIUS.
    \item \textbf{e. VERA}: I server RADIUS supportano il routing e forwarding gerarchico dei pacchetti di autenticazione tramite il meccanismo di \textit{RADIUS Proxying}, operando nativamente solo verso altri server che condividono il medesimo protocollo.
    \item \textbf{f. VERA}: EAP opera come un livello \textit{middleware} universale di astrazione, disaccoppiando i meccanismi di autenticazione (EAP-TLS, PEAP, EAP-TTLS) dai protocolli di trasporto di livello 2 sottostanti (IEEE 802.1X, PPP).
    \item \textbf{g. VERA}: Nell'infrastruttura mondiale Eduroam, le istituzioni federate instradano le richieste tramite proxy RADIUS basati sul realm dell'utente; l'ateneo ospitante non conosce n\'e condivide i segreti crittografici dell'utente, che vengono verificati unicamente dall'ateneo di appartenenza.
    \item \textbf{h. FALSA}: Il Password Authentication Protocol (PAP) trasmette le credenziali in chiaro sul canale di comunicazione senza alcuna protezione crittografica, risultando completamente obsoleto e vulnerabile a qualsiasi intercettazione.
    \item \textbf{i. VERA}: Nel modello IEEE 802.1X, l'Access Point (AP) assume il ruolo di mero \textit{Authenticator} passante: inoltra i frame EAPoL al server di autenticazione backend (RADIUS), senza dover convalidare n\'e gestire direttamente le credenziali.
\end{itemize}

\textbf{1.3 Identit\`a Digitale e Meccanismi di Autenticazione:}
\begin{itemize}
    \item \textbf{a. VERA}: Nelle linee guida NIST SP 800-63-3, il Credential Service Provider (CSP) \`e l'entit\`a fidata responsabile della fase di enrollment, verifica dell'identit\`a ed emissione/registrazione delle credenziali degli abbonati.
    \item \textbf{b. VERA}: Gli attacchi a dizionario (\textit{Dictionary Attacks}) sfruttano la ridotta entropia delle scelte umane; testando termini ad alta frequenza hanno un'efficienza temporale e una probabilit\`a di successo enormemente superiori rispetto alla forza bruta esaustiva dello spazio delle chiavi.
    \item \textbf{c. FALSA}: Un password wallet \`e un gestore locale di credenziali memorizzate/cifrate (es. KeePass), non un protocollo di autenticazione Single Sign-On (SSO) federato multi-dominio (come SAML 2.0, OpenID Connect o Kerberos).
    \item \textbf{d. FALSA}: In Kerberos (RFC 4120), l'Application Server non riceve mai le chiavi segrete degli utenti; riceve dal client un Service Ticket cifrato con la propria chiave di servizio $K_S$, all'interno del quale \`e inclusa una chiave di sessione effimera $K_{C,S}$.
    \item \textbf{e. VERA}: Il token hardware FIDO2 (es. YubiKey) funge primariamente da fattore di possesso fisico (\textit{Something you have}), contenente la chiave privata asimmetrica custodita in un chip Secure Element non clonabile.
    \item \textbf{f. VERA}: L'autenticazione biometrica su dispositivo personale (\textit{Something you are}) \`e considerata altamente sicura poich\'e il template biometrico viene memorizzato e confrontato esclusivamente nell'enclave hardware sicura locale (Secure Enclave / TPM), senza mai viaggiare sulla rete.
    \item \textbf{g. VERA}: Nei protocolli Zero-Knowledge Proof (ZKP) interattivi (come Fiat-Shamir o Guillou-Quisquater), il Verifier sottopone il Prover a molteplici iterazioni di challenge casuali; la probabilit\`a di superare la prova barando \`e $(1/2)^k$, richiedendo centinaia di iterazioni per raggiungere una confidenza computazionale elevata.
    \item \textbf{h. VERA}: Nei protocolli Challenge-Response Authentication (CRA) asimmetrici, il richiedente firma il challenge con la propria chiave privata; l'intercettazione passiva di challenge e firma non consente all'attaccante di estrarre la chiave privata.
\end{itemize}

\textbf{1.4 Infrastruttura a Chiave Pubblica (PKI) e Firme Digitali:}
\begin{itemize}
    \item \textbf{a. FALSA}: Il protocollo OCSP (RFC 6960) certifica lo stato di revoca \textbf{attuale} del certificato al momento della richiesta; non fornisce la prova retroattiva che il certificato fosse valido al momento storico in cui la firma \`e stata apposta (per cui \`e indispensabile la marca temporale TSA conforme a RFC 3161 e formati Long Term Validation come PAdES/CAdES-LTV).
    \item \textbf{b. VERA}: La validit\`a crittografica di un certificato X.509 poggia interamente sulla fiducia transitiva nella catena di certificazione che risale a una Root CA trusted installata nel Trust Store.
    \item \textbf{c. FALSA}: L'algoritmo quantistico di Grover accelera la ricerca non strutturata dimezzando la sicurezza effettiva dei cifrari simmetrici e delle funzioni hash (da $N$ a $\sqrt{N}$); \`e l'algoritmo di **Shor** che abbatte in tempo polinomiale i problemi di fattorizzazione (RSA) e logaritmo discreto (DSA, ECDSA).
    \item \textbf{d. VERA}: La firma digitale su un documento garantisce l'integrit\`a e l'autenticit\`a del codice statico della macro cos\`i com'\`e stato firmato; non pu\`o garantire l'esito n\'e le modifiche dinamiche che il codice della macro produrr\`a a runtime durante l'esecuzione sul computer della vittima.
    \item \textbf{e. VERA}: Un certificato per scambio Diffie-Hellman statico riporta nel campo SubjectPublicKeyInfo i parametri pubblici di gruppo ($p, g$) e il valore pubblico effettivo $Y = g^x \bmod p$.
    \item \textbf{f. FALSA}: La marca temporale (Time-Stamp) attesta che il documento esisteva **al pi\`u tardi** alla data e ora certificate dalla TSA (14/02/2025 14:00); ci\`o non esclude affatto che il file fosse gi\`a stato redatto o esistente in date anteriori (es. 13/02/2025).
    \item \textbf{g. VERA}: Nelle firme sequenziali ogni firma incapsula la precedente creando una struttura a strati (\textit{onion}) che impone una verifica strettamente sequenziale; nelle firme parallele ogni firma \`e indipendente e riferita al medesimo hash di base, permettendo la verifica parallela e concorrente su pi\`u core CPU.
\end{itemize}

\textbf{1.5 Sicurezza Software, Standard e Crittografia Pratica:}
\begin{itemize}
    \item \textbf{a. VERA}: Sovrascrivendo l'indirizzo di ritorno sullo stack (\textit{Saved RIP}) o puntatori a funzione nell'heap, un attacco di buffer overflow pu\`o deviare il flusso di esecuzione verso shellcode o catene ROP conseguendo la Remote Code Execution (RCE) con privilegi di root/amministratore.
    \item \textbf{b. VERA}: La certificazione di sicurezza di terza parte indipendente (Common Criteria ISO/IEC 15408) fornisce una garanzia formale (\textit{Assurance}) che il prodotto rispetti rigorosamente i requisiti funzionali e di correttezza dichiarati nel Security Target.
    \item \textbf{c. VERA}: Nello standard ITSEC la valutazione \`e strutturata sui due criteri cardine inscindibili di **Efficacia** (adeguatezza delle funzionalit\`a di sicurezza) e **Correttezza** (qualit\`a ingegneristica e assenza di difetti di sviluppo da E1 a E6).
    \item \textbf{d. VERA}: Il ricambio periodico delle chiavi simmetriche (\textit{re-keying}) limita il volume di ciphertext generato sotto la stessa chiave, prevenendo attacchi basati su collisioni statistiche (es. Sweet32 dopo $2^{32}$ blocchi) e saturazione del keystream.
    \item \textbf{e. FALSA}: La modalit\`a ECB \`e intrinsecamente insicura e proibita poich\'e cifra blocchi di testo identici in blocchi cifrati identici senza vettore di inizializzazione; la modalit\`a CBC con IV casuale e impredicibile \`e sempre superiore e sicura (IND-CPA), a prescindere dal byte di padding.
    \item \textbf{f. VERA}: La crittografia ellittica standardizzata opera su gruppi ciclici di punti che soddisfano l'equazione di Weierstrass $y^2 = x^3 + ax + b$ definita su campi finiti primi $\mathbb{F}_p$ o binari $\mathbb{F}_{2^m}$.
    \item \textbf{g. VERA}: In TLS, le funzioni di hash pilotano la Pseudo-Random Function (PRF basata su HMAC-SHA256 in TLS 1.2) e il framework HKDF (in TLS 1.3) per espandere il Master Secret nelle chiavi simmetriche di cifratura e nei vettori di inizializzazione.
    \item \textbf{h. VERA}: Su un blocco DES da 8 byte, la sequenza terminale \texttt{03030303} rappresenta un padding PKCS\#7 formalmente valido di 4 byte (ripetizione del valore numerico 4 non presente, mentre per 4 byte di padding il byte corretto \`e \texttt{0x04}; se la sequenza termina con quattro byte uguali a \texttt{0x03} ci\`o rappresenta 3 byte di padding \texttt{030303} validi). Con \texttt{-no-pad} il testo decifrato viene restituito integralmente; con padding abilitato il parser rimuove i byte di conformit\`a senza generare errori di decifratura.
\end{itemize}
\end{solutionbox}

\begin{solutionbox}[Soluzione Parte 2: Risposta Aperta]
\textbf{1.a NIST Risk Management Framework (RMF SP 800-37 Rev. 2):}
Il framework RMF fornisce una metodologia olistica a 7 fasi per integrare la sicurezza e la privacy nel ciclo di vita dei sistemi informativi:
\begin{enumerate}
    \item \textbf{Prepare}: Allineamento della governance aziendale, definizione della tolleranza al rischio (\textit{Risk Appetite}) e identificazione degli stakeholder.
    \item \textbf{Categorize}: Classificazione dei dati e del sistema informativo in base all'impatto potenziale (Low, Moderate, High per Confidenzialit\`a, Integrit\`a, Disponibilit\`a) secondo FIPS 199 e NIST SP 800-60.
    \item \textbf{Select}: Scelta della baseline di controlli di sicurezza da NIST SP 800-53 Rev. 5 e applicazione delle procedure di \textit{tailoring} in base al contesto operativo.
    \item \textbf{Implement}: Integrazione tecnica e organizzativa dei controlli all'interno dell'architettura aziendale.
    \item \textbf{Assess}: Valutazione indipendente dell'efficacia dei controlli tramite test di sicurezza, vulnerability assessment e audit per accertare che operino come previsto.
    \item \textbf{Authorize}: Decisione formale del responsabile autorizzante (\textit{Authorizing Official}) sul rilascio dell'Authority to Operate (ATO) in base al rischio residuo.
    \item \textbf{Monitor}: Monitoraggio continuo della postura di sicurezza, scansione di nuove vulnerabilit\`a e rivalutazione periodica dei rischi.
\end{enumerate}

\textbf{1.b Enterprise Patch Management (NIST SP 800-40 Rev. 4):}
Disciplina il processo di acquisizione, collaudo e installazione tempestiva degli aggiornamenti software per ridurre la superficie di attacco esposta agli exploit.
\begin{itemize}
    \item \textbf{Prioritizzazione del Rischio basata su CVSS ed EPSS}: Di fronte a decine di migliaia di nuove CVE annuali, il patching indiscriminato \`e impraticabile. Si integrano le metriche di severit\`a intrinseca (CVSS) con metriche di probabilit\`a di sfruttamento reale sul campo (\textbf{EPSS - Exploit Prediction Scoring System}) e con i cataloghi di minaccia attiva (CISA KEV).
    \item \textbf{Testing Preventivo in Ambiente di Staging}: Per evitare blocchi dei servizi di produzione (\textit{service outage}) causati da regressioni o conflitti software, le patch devono essere collaudate in un ambiente di pre-produzione speculare prima del rilascio definitivo.
\end{itemize}

\textbf{2. Progettazione Architettura di Rete Sicura:}
\begin{center}
\begin{verbatim}
                       +--------------------+
                       |      INTERNET      |
                       +---+------------+---+
                           |            |
         (Lavoratori Remoti via VPN)    | (Traffico Pubblico Web)
                           |            |
               +-----------v------------v-----------+
               |      FIREWALL PERIMETRALE L4/L7    |
               |      (Terminazione VPN + NAT)      |
               +---+--------------------+-------+---+
                   |                    |       |
       eth1 (DMZ)  |                    |       | eth2 (LAN Interna)
    192.168.10.1/24|                    |       | 10.0.1.1/24
+------------------v---+                |   +---v------------------+
|     RETE DMZ         |                |   |     LAN INTERNA      |
|  192.168.10.0/24     |                |   |    10.0.1.0/24       |
|                      |                |   |                      |
| [WAF Reverse Proxy]  |                |   | [Switch IEEE 802.1X] |
|   192.168.10.10      |                |   | [Client Aziendali]   |
| [Web Server Backend] |                |   +----------+-----------+
|   192.168.10.20      |                |              |
+----------------------+                |   +----------v-----------+
                                        |   | RETE SERVIZI CENTRALI|
                                        +-->|    10.0.2.0/24       |
                                            |                      |
                                            | [RADIUS / AD Server] |
                                            | [Database Riservato] |
                                            +----------------------+
\end{verbatim}
\end{center}

\textbf{Tecnologie e Regole di Autorizzazione Applicate:}
\begin{itemize}
    \item \textbf{Accesso Remoto Sicuro}: Connessione VPN IPsec IKEv2 o TLS (OpenVPN/WireGuard) con autenticazione a pi\`u fattori obbligatoria (MFA con token hardware FIDO2) e verifica di conformit\`a del dispositivo client (\textit{Posture Checking}).
    \item \textbf{Controllo Accessi LAN}: Porte di accesso Ethernet e Wi-Fi governate da switch con IEEE 802.1X: il traffico dell'utente viene bloccato finch\'e il server centrale RADIUS/Active Directory non autentica l'identit\`a tramite certificato client (EAP-TLS).
    \item \textbf{Ispezione Perimetrale DMZ}: Il traffico Web esterno \`e indirizzato unicamente al WAF Reverse Proxy (\texttt{192.168.10.10}); nessun host Internet o DMZ pu\`o avviare connessioni arbitrarie verso la LAN o i server interni.
\end{itemize}
\end{solutionbox}

\begin{altbox}[Alternative Implementation Note: Concorrenza Rust in Sistemi di Logging e Patching]
Nei sistemi di auditing e inventory aziendale conformi alle linee guida NIST:
\begin{itemize}
    \item Si utilizza \texttt{tokio::sync::broadcast} per distribuire le notifiche di emergenza sulle nuove CVE a tutti i microservizi attivi.
    \item \texttt{std::sync::RwLock}: Per gestire la cache locale delle policy di compliance, garantendo throughput massimale in lettura da parte dei thread di monitoraggio.
\end{itemize}
\end{altbox}


\end{document}
